%% ma: L10-B15-0051
%% lop: 10
%% chuong: 6
%% bai: 15
%% dang: 10.15.6
%% loai: TL
%% muc: TH
%% dap_an: "-"
%% nguon: "Ảnh giáo viên gửi 21/11/2026: Chuyên đề VI – Toán 10 – Chương VI (Hàm số – đồ thị và ứng dụng), tr. 199–203, Dạng 6 Câu 1"
%% kien_thuc: "Lớp 10 Bài 15 (hàm số trong thực tế)"
%% trang_thai: chinh-thuc
%% ghi_chu: "Dạng do giáo viên cung cấp (21/11/2026)."
\begin{ex}
Cho rằng diện tích rừng nhiệt đới trên trái đất được xác định bởi hàm số $S=718{,}3-4{,}6t$, trong đó $S$ được tính bằng triệu hec-ta, $t$ tính bằng số năm kể từ năm $1990$. Hãy tính diện tích rừng nhiệt đới vào các năm $1990$ và $2018$.
\loigiai{Năm $1990$: $t=0$, $S=718{,}3$ triệu ha. Năm $2018$: $t=28$, $S=718{,}3-4{,}6\cdot28=589{,}5$ triệu ha.}
\end{ex}
